Une petite fusée de \qty{\v{M}}{kg} décolle verticalement. Ses moteurs exercent une poussée constante de \qty{\v{F}}{kN} pendant \qty{\v{t}}{s}, puis s'arrêtent. On suppose que la masse reste constante et on néglige la résistance de l'air.

\begin{enumerate}[label=\alph*)]
    \item Faites le diagramme des forces de la fusée pendant la poussée et déterminez son accélération.
    \item Quelle est son altitude et sa vitesse lorsque les moteurs s'arrêtent ?
    \item Quelle altitude maximale la fusée atteint-elle ?
\end{enumerate}
